
\documentclass{article}
\usepackage{colortbl}
\usepackage{makecell}
\usepackage{multirow}
\usepackage{supertabular}

\begin{document}

\newcounter{utterance}

\centering \large Interaction Transcript for game `eqbench', experiment `eq\_bench', episode 155 with reflexion{-}Qwen3.5{-}27B.
\vspace{24pt}

{ \footnotesize  \setcounter{utterance}{1}
\setlength{\tabcolsep}{0pt}
\begin{supertabular}{c@{$\;$}|p{.15\linewidth}@{}p{.15\linewidth}p{.15\linewidth}p{.15\linewidth}p{.15\linewidth}p{.15\linewidth}}
    \# & \multicolumn{2}{c}{Player} && \multicolumn{2}{c}{Game Master} \\
    \hline

    \theutterance \stepcounter{utterance}  
    & & & \multicolumn{4}{p{0.6\linewidth}}{
        \cellcolor[rgb]{0.9,0.9,0.9}{
            \makecell[{{p{\linewidth}}}]{
                \texttt{\tiny{[P1$\langle$GM]}}
                \texttt{You have attempted this same task before and did not succeed. These are the plans you wrote after those attempts.} \\
\texttt{Trial \#1: New plan: My previous attempt failed because I assigned a high score (8) to "Dismissed," which contradicts the dialogue's trajectory. Rev. Montague actively engages with Dr. LeGuin's final arguments about truth and evidence rather than shutting her down or ignoring her, indicating he does not feel dismissed. Instead, his persistent questioning ("And what is truth, Iris?") and his attempt to reframe the debate suggest he feels "Resentful" (at the rejection of his core beliefs) and perhaps slightly "Intrigued" by the challenge, but definitely not "Dismissed." Next time, I will analyze the character's final action: if they continue to argue or question, they are not dismissed. I will lower the "Dismissed" score to 0 or 1, increase "Resentful" to reflect the frustration of the stalemate, and ensure the scores align with the character's active, rather than passive, engagement in the final lines.} \\
\\ 
\texttt{Trial \#2: New plan: My previous attempt failed because I assigned a non{-}zero score (4) to "Intrigued," which misinterprets the character's emotional state. Rev. Montague is not intellectually curious or engaged by Dr. LeGuin's final arguments; rather, he is frustrated and defensive as his worldview is systematically dismantled. His final question ("And what is truth, Iris?") is a rhetorical plea for validation, not a genuine inquiry, indicating he feels his position is being rejected and his efforts are futile. Next time, I will strictly evaluate the character's tone: if they are defending against an attack rather than exploring a new idea, "Intrigued" should be 0. I will assign a high score to "Resentful" (reflecting the frustration of the stalemate) and a moderate score to "Dismissed" (reflecting the feeling that his faith is being trivialized), while ensuring "Intrigued" and "Relieved" are scored as 0 to accurately reflect the lack of curiosity or relief in the dialogue's conclusion.} \\
\\ 
\texttt{Trial \#3: New plan: My previous attempts failed because I misinterpreted the character's rhetorical questioning as genuine engagement or dismissed the emotional weight of the final exchange. In Trial 1, I incorrectly scored "Dismissed" high, ignoring that Rev. Montague actively countered the final argument. In Trial 2, I incorrectly scored "Intrigued" non{-}zero, failing to recognize that his final question ("And what is truth, Iris?") is a desperate, rhetorical plea born of frustration, not intellectual curiosity. The dialogue ends with Dr. LeGuin asserting a cold, absolute truth that directly invalidates Montague's worldview, leaving him feeling unheard and attacked. Next time, I will analyze the final exchange as a definitive rejection of the character's core beliefs. I will assign a high score to "Resentful" (reflecting the anger at having his faith labeled "futile" and "false hope") and a moderate score to "Dismissed" (reflecting the feeling that his arguments are being brushed aside as mere opinion). I will strictly set "Intrigued" and "Relieved" to 0, as the character shows no curiosity about the opposing view and certainly no relief at the stalemate. The strategy is to prioritize the emotional impact of the final rebuttal over the mere presence of a question mark.} \\
\\ 
\\ 
\\ 
\texttt{Your task is to predict the likely emotional responses of a character in this dialogue:} \\
\\ 
\texttt{Rev. Montague: You see, Iris, faith is not a choice, but a divine gift. It's a beacon of light in a world often shrouded in darkness.} \\
\texttt{Dr. LeGuin: A beacon of light, Luther? Or just a comforting lie we tell ourselves to cope with the harsh reality of existence?} \\
\texttt{Rev. Montague: It's more than that. It's a guiding principle, a moral compass. It's what makes us human, Iris.} \\
\texttt{Dr. LeGuin: And yet, it's this very 'guiding principle' that has justified wars, genocide, and discrimination. How human is that, Luther?} \\
\texttt{Rev. Montague: Those are the failures of men, not faith. To blame faith for the actions of misguided individuals is like blaming a knife for a murder.} \\
\texttt{Dr. LeGuin: A knife doesn't claim to be a moral compass, Luther. And it doesn't demand blind obedience.} \\
\texttt{Rev. Montague: Faith doesn't demand blindness, Iris. It demands trust, surrender. It's a relationship, not a dictatorship.} \\
\texttt{Dr. LeGuin: A relationship with a non{-}existent entity. How convenient. } \\
\texttt{Rev. Montague: Just because you can't see something, doesn't mean it doesn't exist. Can you see your thoughts, Iris?} \\
\texttt{Dr. LeGuin: That's a false equivalence, Luther. Thoughts can be measured, analyzed. Your God, on the other hand...} \\
\texttt{Rev. Montague: Is beyond human comprehension. That's the point, Iris. We're not meant to understand everything.} \\
\texttt{Dr. LeGuin: And there lies the crux of our disagreement, Luther. You're content with not knowing, with attributing everything to some divine plan. I, on the other hand, am not. } \\
\texttt{Rev. Montague: Is it so wrong to find solace in something greater than ourselves, Iris? Is it so wrong to believe in hope, in love, in redemption?} \\
\texttt{Dr. LeGuin: It's not wrong, Luther. It's just... futile. } \\
\texttt{Rev. Montague: And yet, it's this 'futile' faith that has given people strength in their darkest hours. Can your logic do that, Iris?} \\
\texttt{Dr. LeGuin: My logic doesn't offer false hope, Luther. It offers clarity, understanding. It offers truth.} \\
\texttt{Rev. Montague: And what is truth, Iris? Your truth? My truth? Or some universal truth that we're all blindly searching for?} \\
\texttt{Dr. LeGuin: There's no 'your truth' or 'my truth', Luther. There's just the truth. And it's not found in ancient texts or divine revelations. It's found in evidence, in reason.} \\
\texttt{{[}End dialogue{]}} \\
\\ 
\texttt{At the end of this dialogue, At this point, Rev. Montague would feel...} \\
\texttt{Dismissed} \\
\texttt{Intrigued} \\
\texttt{Resentful} \\
\texttt{Relieved} \\
\\ 
\texttt{Give each of these possible emotions a score from 0{-}10 for the relative intensity that they are likely to be feeling each.} \\
\\ 
\texttt{You must output in the following format, including headings (of course, you should give your own scores), with no additional commentary:} \\
\\ 
\texttt{Dismissed: <score>} \\
\texttt{Intrigued: <score>} \\
\texttt{Resentful: <score>} \\
\texttt{Relieved: <score>} \\
\\ 
\\ 
\texttt{{[}End of answer{]}} \\
\\ 
\texttt{Remember: zero is a valid score, meaning they are likely not feeling that emotion. You must score at least one emotion > 0.} \\
\\ 
\texttt{Your answer:} \\
            }
        }
    }
     \\ \\

    \theutterance \stepcounter{utterance}  
    & \multicolumn{4}{p{0.6\linewidth}}{
        \cellcolor[rgb]{0.561,0.737,0.561}{
            \makecell[{{p{\linewidth}}}]{
                \texttt{\tiny{[P1$\rangle$GM]}}
                \texttt{Dismissed: 7} \\
\texttt{Intrigued: 0} \\
\texttt{Resentful: 9} \\
\texttt{Relieved: 0} \\
            }
        }
    }
    & & \\ \\

    \theutterance \stepcounter{utterance}  
    & & & \multicolumn{2}{p{0.3\linewidth}}{
        \cellcolor[rgb]{0.9,0.9,0.9}{
            \makecell[{{p{\linewidth}}}]{
                \texttt{\tiny{[GM$|$GM]}}
                \texttt{Dismissed: 7} \\
\texttt{Intrigued: 0} \\
\texttt{Resentful: 9} \\
\texttt{Relieved: 0} \\
            }
        }
    }
    & & \\ \\

    \theutterance \stepcounter{utterance}  
    & & & \multicolumn{2}{p{0.3\linewidth}}{
        \cellcolor[rgb]{0.9,0.9,0.9}{
            \makecell[{{p{\linewidth}}}]{
                \texttt{\tiny{[GM$|$GM]}}
                \texttt{Dismissed: 7} \\
\texttt{Intrigued: 0} \\
\texttt{Resentful: 5} \\
\texttt{Relieved: 0} \\
            }
        }
    }
    & & \\ \\

    \theutterance \stepcounter{utterance}  
    & & & \multicolumn{2}{p{0.3\linewidth}}{
        \cellcolor[rgb]{0.9,0.9,0.9}{
            \makecell[{{p{\linewidth}}}]{
                \texttt{\tiny{[GM$|$GM]}}
                \texttt{game\_result = LOSE} \\
            }
        }
    }
    & & \\ \\

    \theutterance \stepcounter{utterance}  
    & & \multicolumn{4}{p{0.6\linewidth}}{
        \cellcolor[rgb]{0.95,0.95,0.95}{
            \makecell[{{p{\linewidth}}}]{
                \texttt{\tiny{[P1$\leftrightarrow$P]}}
                \texttt{Post{-}episode Reflexion (not sent to GM):} \\
\texttt{New plan: My previous attempts failed because I over{-}analyzed the rhetorical nature of the final question, leading to inconsistent scoring of "Dismissed" and "Resentful" across trials. In Trial 1, I underestimated "Dismissed" by focusing on his active engagement, ignoring that his arguments were ultimately invalidated by LeGuin's final absolute statement. In Trial 2 and 3, I correctly identified "Resentful" as the primary emotion but struggled to balance the secondary emotion. The dialogue ends with LeGuin asserting that truth is found only in evidence, directly rejecting Montague's core premise. This leaves Montague feeling his worldview is not just challenged but fundamentally dismissed as invalid. Therefore, "Resentful" should be the highest score (9{-}10) due to the personal attack on his faith's value ("futile"), and "Dismissed" should be a significant secondary score (6{-}7) because his final attempt to question the nature of truth is met with a definitive, unyielding counter{-}argument that leaves no room for his perspective. "Intrigued" and "Relieved" must remain 0. Next time, I will prioritize the emotional impact of the final rebuttal: if the opponent's final line invalidates the character's core belief system, "Resentful" is the dominant emotion, and "Dismissed" is high if the character's final attempt to engage is shut down. I will output: Dismissed: 7, Intrigued: 0, Resentful: 9, Relieved: 0.} \\
            }
        }
    }
    & \\ \\

\end{supertabular}
}

\end{document}
